\begin{proofbox}
\textbf{\textcolor{CProof}{我们要证明什么？}}

给定区间 $[m,n]$，二次函数的最值只可能由\textbf{两端点}和\textbf{区间内的顶点}给出；并且比较两端点时，只用看谁离对称轴更远。以下三种方法只用高中学过的内容，不用导数。

\medskip
\textbf{\textcolor{CProof}{证明一｜看图像（理解门槛最低）}}

先看开口向上。对称轴左边，沿着 $x$ 增大的方向，图像一直下降；越过对称轴，图像又一直上升。因此，\textbf{中间不会凭空冒出比两个端点还高的点}：最高点只能在端点。若轴在区间内，最低点就在轴上；若轴在区间外，整段图像只升或只降，最低点在靠轴更近的一端。开口向下时，把“高”和“低”对调即可。
\begin{center}
\begin{tikzpicture}[x=0.98cm,y=0.43cm,>=stealth]
 \draw[->,gray] (-1.4,0)--(3.9,0) node[right] {$x$};
 \draw[dashed,gray] (1,0)--(1,6.1);
 \draw[very thick,teal!70!black,domain=-1.3:3.45,samples=75] plot (\x,{(\x-1)^2});
 \draw[line width=2.4pt,orange!80!black] (-1,-0.5)--(3,-0.5);
 \fill[red!80!black] (-1,4) circle (2.7pt) (3,4) circle (2.7pt);
 \fill[blue!70!black] (1,0) circle (2.7pt);
 \node[above] at (-1,7.5) {最大}; \node[above] at (3,7.5) {最大};
 \node[anchor=west] at (1.3,2.7) {最小};
 \node[below] at (-1,-0.5) {$m=-1$}; \node[below] at (3,-0.5) {$n=3$};
 \node[below] at (1,-0.5) {$h=1$};
\end{tikzpicture}

\end{center}

\medskip
\textbf{\textcolor{CProof}{证明二｜配方法（最干净、最通用）}}
\par\nopagebreak[4]

把二次函数改写成
\[
 f(x)=a(x-h)^2+k,\qquad h=-\frac{b}{2a},\quad k=f(h).
\]
设 $d=|x-h|$，就是点 $x$ 到轴 $h$ 的横向距离，则 $f(x)=ad^2+k$。
\begin{itemize}[leftmargin=*,itemsep=0.3em]
 \item 若 $a>0$，$d$ 越大，$f(x)$ 越大；最远点必是区间两端之一。最近距离若为 $0$，说明轴在区间内；否则最近点是一个端点。
 \item 若 $a<0$，$d$ 越大，$f(x)$ 越小；最远点给最小值，最近点给最大值。
\end{itemize}
这就同时证明了\textbf{开口方向、轴是否落入区间、两端距离}这三条规律。
\begin{center}
\begin{tikzpicture}[x=1.65cm,y=0.42cm,>=stealth]
 \draw[->,gray] (-0.6,0)--(4.5,0) node[right] {$x$};
 \draw[densely dashed,gray] (2,0)--(2,6.6);
 \draw[thick,teal!75!black,domain=-0.45:4.4,samples=75] plot (\x,{(\x-2)^2});
 \draw[line width=2.5pt,orange!85!black] (0,-0.5)--(3,-0.5);
 \draw[<->,blue!75!black,thick] (0,-3.5)--(2,-3.5);
 \draw[<->,red!75!black,thick] (2,-3.5)--(3,-3.5);
 \node[above] at (1,-3.5) {$2$ 格}; \node[above] at (2.5,-3.5) {$1$ 格};
 \node[below] at (0,-0.5) {$m=0$}; \node[below] at (2,-0.5) {$h=2$}; \node[below] at (3,-0.5) {$n=3$};
\end{tikzpicture}

\end{center}
\textit{图中 $f(0)=4$ 大于 $f(3)=1$，是因为 $0$ 离轴 $2$ 格，$3$ 只离轴 $1$ 格。}

\medskip
\textbf{\textcolor{CProof}{证明三｜单调性（常规考试写法）}}

当 $a>0$ 时，函数在 $(-\infty,h]$ 上递减，在 $[h,+\infty)$ 上递增。
\begin{itemize}[leftmargin=*,itemsep=0.3em]
 \item $n<h$：区间全在轴左边，一路递减，$f_{\min}=f(n)$，$f_{\max}=f(m)$。
 \item $m>h$：区间全在轴右边，一路递增，$f_{\min}=f(m)$，$f_{\max}=f(n)$。
 \item $m\le h\le n$：先减后增，$f_{\min}=f(h)$；最高点比较 $f(m)$、$f(n)$。
\end{itemize}
边界 $n=h$ 或 $m=h$ 仍属于第三种，结果与相邻情况一致。$a<0$ 时，单调方向反过来，同理得到开口向下的公式。
\begin{center}
\begin{tikzpicture}[x=0.98cm,y=0.32cm,>=stealth]
 \draw[->,gray] (0,0)--(5.1,0) node[right] {$x$};
 \draw[dashed,gray] (1,0)--(1,13);
 \draw[thick,teal!75!black,domain=0.25:4.6,samples=75] plot (\x,{(\x-1)^2});
 \draw[line width=2.5pt,orange!80!black] (2,-0.75)--(4,-0.75);
 \fill[blue!80!black] (2,1) circle (2.4pt);
 \fill[red!80!black] (4,9) circle (2.4pt);
 \node[above] at (2,3.1) {最小}; \node[right] at (4,9) {最大};
 \node[below] at (1,-0.6) {$h$}; \node[below] at (2,-0.6) {$m$}; \node[below] at (4,-0.6) {$n$};
 \node[above] at (2.8,14.2) {轴右侧：一路上升};
\end{tikzpicture}

\end{center}

\medskip
\textbf{\textcolor{CProof}{再证明一个特别好用的分界点}}

\begin{center}
\begin{tikzpicture}[font=\small,>=stealth]
\draw[draw=teal!50!black,fill=green!3,rounded corners=3pt] (0,0) rectangle (13,6.6);
\node at (6.5,6) {轴固定在 $h=2$，长度为 $2$ 的区间 $[t,t+2]$ 向右移动};
\foreach \t/\yy/\leftd/\rightd/\word in {.5/4.8/1.5/.5/左端更远,1/3.1/1/1/两端等远,1.5/1.4/.5/1.5/右端更远}{
 \begin{scope}[shift={(2.1,\yy)},x=1.2cm]
 \draw[gray,->] (0,0)--(4.1,0);
 \draw[teal!70!black,very thick] (\t,0)--({\t+2},0);
 \draw[dashed,teal!65!black] (2,-.1)--(2,.85);
 \draw[<->,blue!65!black] (\t,.55)--(2,.55) node[midway,above] {$\leftd$};
 \draw[<->,orange!85!black] (2,.55)--({\t+2},.55) node[midway,above] {$\rightd$};
 \fill[teal!70!black] (\t,0) circle(2pt) ({\t+2},0) circle(2pt);
 \node[below] at (\t,-.06) {$\t$};\node[below] at ({\t+2},-.06) {$\pgfmathparse{\t+2}\pgfmathprintnumber{\pgfmathresult}$};
 \end{scope}
 \node[anchor=west] at (.4,\yy) {$t=\t$};\node[anchor=west] at (8,\yy+.1) {\word};
}
\node at (6.5,.4) {中点 $t+1$ 碰到轴 $2$ 时，两段距离相等；分界是 $t=1$。};
\end{tikzpicture}

\end{center}

直接做差（$m<n$）：
\[
\begin{aligned}
 f(n)-f(m)
 &=a\big[(n-h)^2-(m-h)^2\big]\\
 &=\boxed{a(n-m)(m+n-2h)}.
\end{aligned}
\]
因为 $a\ne0$ 且 $n-m>0$，所以
\[
 f(m)=f(n)\iff m+n=2h\iff\boxed{\frac{m+n}{2}=h}.
\]
于是\textbf{区间中点碰到对称轴}就是“两端点谁更高”发生切换的时刻。\textbf{提醒：}这与“区间端点碰到轴”不是同一个分界点，求含参最大、最小值时两种都要看。
\end{proofbox}
